\documentclass[11pt]{article}

\usepackage[margin=1in]{geometry}
\usepackage[T1]{fontenc}
\usepackage[utf8]{inputenc}
\usepackage{amsmath,amssymb,mathtools}
\usepackage{array}
\usepackage{booktabs}
\usepackage{enumitem}
\usepackage{graphicx}
\usepackage{longtable}
\usepackage[expansion=false]{microtype}
\usepackage{hyperref}
\usepackage[numbers,sort&compress]{natbib}
\usepackage{cleveref}
\emergencystretch=2em

\newcolumntype{L}[1]{>{\raggedright\arraybackslash\sloppy}p{#1}}
\setlength{\tabcolsep}{5pt}
\renewcommand{\arraystretch}{1.12}

\hypersetup{
  colorlinks=true,
  linkcolor=blue,
  citecolor=blue,
  urlcolor=blue
}

\newcommand{\E}{\mathbb{E}}
\newcommand{\NPV}{\mathrm{NPV}}
\newcommand{\code}[1]{\nolinkurl{#1}}

\title{Gap-Closing Literature Note for the Credit-Card Customer NPV Proposal\\
What Existing Literatures Can and Cannot Close}
\author{Chak Wong}
\date{July 2, 2026}

\begin{document}
\maketitle

\begin{abstract}
This note is a companion to the credit-card customer NPV proposal. Its purpose
is not to restate the proposal or to replace the stochastic-process note. It is
to map the proposal's remaining theoretical and modeling gaps to the relevant
literatures and to state, gap by gap, what those literatures can actually
support. The note is therefore a committee-facing crosswalk. For each major gap,
it identifies the exact proposal sections and equations where the issue matters,
explains why the gap is material, cites the literature families that can close
or tighten it, distinguishes literature support from proposal derivation and
bank-specific empirical work, and states what remains unresolved even after the
literature is added. The note is intentionally conservative: many gaps can be
narrowed substantially by existing literature, but few are eliminated without
issuer-specific data, experiments, and governance choices.
\end{abstract}

\section{Introduction and scope}

The current proposal has become much stronger on valuation semantics,
component decomposition, and dynamic structure. It now contains a long section
that assembles the proposal into a Bayesian state-space / belief-state control
system. Even so, a committee of econometrics and stochastic-control reviewers
can still press on deeper theoretical and modeling gaps. Some are about
mathematical structure: is the state sufficient, is the filtration explicit, is
the joint kernel coherent? Others are about modeling and identification: how is
outside wallet identified, what is the state unit, when is relationship value
causal rather than confounded, and where exactly does fixed-policy valuation end
and optimization begin?

This note answers those objections in a structured way. It does not claim that
literature alone solves the bank's problem. Instead, it uses literature to do
four narrower jobs:
\begin{enumerate}[label=\arabic*.]
  \item justify formal language or abstraction,
  \item justify a modeling move or computational approximation,
  \item show what has already been studied in nearby dynamic decision settings,
  \item and mark the line between what can be defended theoretically and what
        still requires internal bank evidence.
\end{enumerate}

\section{How to read this note}

Each gap section follows the same template.
\begin{enumerate}[label=\arabic*.]
  \item \textbf{Gap in current proposal.} What remains under-specified or
        vulnerable.
  \item \textbf{Where the gap appears.} Exact section or equation references in
        the main proposal.
  \item \textbf{Literatures that can close or tighten it.} The literature family
        and key citations.
  \item \textbf{What the literature actually supports.} The strongest claim that
        can be defended.
  \item \textbf{What remains bank-specific.} The part that still must be treated
        as an issuer-specific empirical object, assumption, or governance rule.
\end{enumerate}

The note uses four support labels throughout:
\begin{enumerate}[label=\arabic*.]
  \item \emph{literature-supported mechanism},
  \item \emph{literature-supported formalism},
  \item \emph{proposal derivation},
  \item \emph{bank-specific empirical object or governance convention}.
\end{enumerate}
A reader should not infer that all gaps can be closed to the same degree.
Several can be tightened greatly at the theoretical level while remaining only
partially resolved empirically.

\section{Crosswalk table}

Proposal equations live in the main proposal file, not in this companion note, so
all equation anchors below are written as plain labels in parentheses rather than
as live \verb|\eqref| cross-references. The purpose is provenance, not internal
LaTeX navigation.

\begin{longtable}{L{0.18\linewidth}L{0.22\linewidth}L{0.22\linewidth}L{0.18\linewidth}L{0.12\linewidth}}
\caption{Gap-to-literature crosswalk for the main proposal.}
\label{tab:gap-crosswalk}\\
\toprule
Gap category & Main proposal anchors & Literature family & What the literature can support & What remains bank-specific \\
\midrule
\endfirsthead
\toprule
Gap category & Main proposal anchors & Literature family & What the literature can support & What remains bank-specific \\
\midrule
\endhead
State unit / aggregation
  & proposal objective (eq:proposal-objective), state kernel (eq:state-kernel-panel), balance stock-flow (eq:balance-stock-flow)
  & state aggregation, approximate DP, CLV state models
  & disciplined aggregation and compression of dynamic states
  & exact bank unit and sufficient aggregation \\
Markov sufficiency / filtration
  & dynamic-policy system (eq:dynamic-policy-system), policy value (eq:policy-value), state-space assembly section
  & MDP, DTR, POMDP, Bayesian filtering
  & explicit information-state and belief-state language
  & whether the chosen bank features are empirically sufficient \\
Shared macro common state
  & category spend (eq:category-spend), PD transition (eq:pd-transition), term/loss mapping (eq:terms-loss)
  & factor models, macro state-space models, scenario design
  & low-dimensional common-shock macro block
  & exact latent macro law and elasticities \\
Joint kernel / factorization
  & state kernel (eq:state-kernel-panel), balance stock-flow (eq:balance-stock-flow), PD transition (eq:pd-transition)
  & dynamic state-space decomposition, modular conditional laws
  & coherent structured approximation to the joint kernel
  & exact factorization validity \\
Belief approximation / compression
  & stochastic-process section belief-state equations
  & Bayesian filtering, belief compression, finite-memory approximations
  & approximate posterior summaries as implementation objects
  & exact compressed belief sufficient for this bank \\
Incremental baseline semantics
  & proposal objective (eq:proposal-objective), NPV functional (eq:panel-npv-functional), policy value (eq:policy-value)
  & causal inference, DTRs, policy evaluation
  & fixed-policy value vs optimized-policy value distinction
  & approved baseline and policy-bundle semantics \\
Wallet / latent demand identification
  & wallet accounting through cannibalization-to-NPV (eq:wallet-accounting)--(eq:cannibalization-to-npv)
  & payment choice, demand/substitution, uplift, partial identification
  & latent wallet / substitution language and bounded identification
  & issuer-specific wallet capture and movable fraction \\
Relationship value integration
  & relationship value (eq:relationship-value), valuation semantics sections
  & CLV, customer equity, banking profitability
  & cautious separation of core card value and spillover value
  & causal spillover magnitude and inclusion rule \\
Valuation vs optimization boundary
  & dynamic-policy system (eq:dynamic-policy-system), policy value (eq:policy-value), allowed-use and governance sections
  & MDP, DTR, OPE, safe improvement, offline RL
  & fixed-policy simulation first, challenger evaluation second
  & production optimization threshold and governance gate \\
\bottomrule
\end{longtable}

\section{State unit and aggregation}

\paragraph{Gap in current proposal.}
The proposal allows the unit $i$ to be a customer, account, prospect,
household, application, or account-action unit, while the dynamic sections also
use wallet, relationship, and all-bank incremental value objects. This creates a
structural question: what is the actual dynamic state unit? If the state is too
fine, the process becomes intractable; if it is too coarse, internal
cannibalization and relationship effects may be misrepresented.

\paragraph{Where the gap appears.}
This issue enters through proposal objective (eq:proposal-objective),
state kernel (eq:state-kernel-panel), wallet accounting (eq:wallet-accounting),
balance stock-flow (eq:balance-stock-flow), and the new state-space assembly
section where observed and latent state blocks are introduced.

\paragraph{Literatures that can close or tighten it.}
State aggregation and approximate dynamic programming literature directly helps
here. The most useful starting point is \citet{bean1987aggregation}, which is a
general operations-research reference for aggregation in dynamic programming.
For a banking or customer-relationship context, CLV and duration/state models
already present in the proposal's bibliography also help define economically
meaningful state blocks \citep{schmittlein1987counting,fader2005bgnbd,
bolton1998duration}. Banking-specific dynamic optimization references such as
\citet{thomas2011creditcardmdp} and \citet{lewis2005relationshippricing} can be
used to show how dynamic pricing/control papers work with deliberately chosen
state abstractions.

\paragraph{What the literature actually supports.}
The literature supports the claim that state aggregation is a legitimate and
often necessary modeling device, and that economically distinct blocks should
not be collapsed carelessly if they carry different transition or reward
semantics. It also supports the use of lifecycle, balance, payment, risk, and
relationship states as meaningful dynamic components.

\paragraph{What remains bank-specific.}
The literature does not choose the bank's unique correct state unit. It does not
prove whether the correct object is account-level, customer-level,
household-level, or a multi-account relationship state. That remains a proposal
and governance choice to be defended against the bank's actual use cases and
internal cannibalization structure.

\section{Markov sufficiency and filtration}

\paragraph{Gap in current proposal.}
The proposal now distinguishes observed state, latent state, and belief state,
but reviewers can still ask whether the chosen information state is truly
sufficient. If omitted history matters beyond the proposal's declared observed
state and posterior state, then the Bellman recursion is not properly grounded.

\paragraph{Where the gap appears.}
This issue appears in the proposal's dynamic-policy language around
dynamic-policy system (eq:dynamic-policy-system) and policy value
(eq:policy-value), and more sharply in the full stochastic-process section where
the filtration and belief-state objects are introduced.

\paragraph{Literatures that can close or tighten it.}
This is one of the clearest places where the literature can add theoretical
weight. The MDP text of \citet{puterman1994mdp}, the dynamic treatment-regime
reference \citet{murphy2003dynamic}, and especially the POMDP / state-space
literature \citep{smallwood1973pomdp,kaelbling1998pomdp,
ledoux2005recursivefilters,durbinkoopman2012ssm,westharrison1997bfdm} provide
the correct language for explicit filtration, belief-state sufficiency, and
recursive updating.

\paragraph{What the literature actually supports.}
These references strongly support the need to write down:
- a pre-action information set,
- a post-observation information set,
- a belief or filtering distribution over latent state,
- and a policy measurable with respect to the available information.
They justify the proposal's move from static feature stacks to an information-
state or belief-state formulation.

\paragraph{What remains bank-specific.}
The literature does not prove that the proposal's chosen observed state is
actually Markov-sufficient for this bank's problem. That depends on the bank's
real data lags, omitted state variables, missingness, delayed collections
signals, future-policy dependence, and empirical adequacy of the latent-state
model.

\section{Shared macro common-shock state}

\paragraph{Gap in current proposal.}
The proposal repeatedly uses macro and labor-market variables in spend, payment,
loss, and scenario language, but a committee can still ask: how exactly should a
common macro state be represented? If every module conditions on macro
variables separately, the model risks being a loose collection of reduced-form
controls rather than one coherent common-shock dynamic system.

\paragraph{Where the gap appears.}
This matters in the macro section, especially category spend
(eq:category-spend), and in the balance, payment, and loss equations such as PD
transition (eq:pd-transition) and term/loss mapping (eq:terms-loss). It also
matters in the state-space assembly section's macro block.

\paragraph{Literatures that can close or tighten it.}
The most useful gap-closing literatures here are dynamic factor and data-rich
macro-state literatures. \citet{bernankeboivineliasz2005favar} and
\citet{bernankeboivin2003datarich} support reducing a high-dimensional macro
environment to a lower-dimensional latent common state, while
\citet{stockwatson2005factorvar} supports the use of dynamic factor structures
as inputs to multivariate dynamic systems.

\paragraph{What the literature actually supports.}
These references support the proposal's use of a low-dimensional shared macro
state or latent factor block that can feed multiple modules jointly---spend,
payment, delinquency, funding, and attrition---rather than treating macro state
as unrelated controls in each equation.

\paragraph{What remains bank-specific.}
The literature does not identify the bank's exact latent macro regime, how many
factors to retain, how geography enters, or the issuer-specific elasticities of
spend, payment, and risk with respect to that macro state.

\section{Joint kernel and factorization}

\paragraph{Gap in current proposal.}
The proposal has many useful module equations, but the committee can still ask
whether they define one coherent joint stochastic kernel or merely a set of local
conditional blocks that may not compose cleanly.

\paragraph{Where the gap appears.}
This issue sits behind state kernel (eq:state-kernel-panel),
balance stock-flow (eq:balance-stock-flow), payment component
(eq:payment-component), line component (eq:line-component), and PD transition
(eq:pd-transition), and it becomes sharper in the state-space assembly where
routing, balances, payments, delinquency, promotions, and relationship value all
appear as blocks in one process.

\paragraph{Literatures that can close or tighten it.}
The state-space and recursive-filtering literature already in the bibliography is
helpful here because it supports writing joint dynamic systems through
structured conditional laws \citep{durbinkoopman2012ssm,westharrison1997bfdm,
doucet2001smc,ledoux2005recursivefilters}. This does not by itself prove the
proposal's factorization, but it supports the idea that modular conditional
specifications are an acceptable way to approximate high-dimensional joint
dynamics.

\paragraph{What the literature actually supports.}
The literature supports structured decomposition of a high-dimensional process
into conditional blocks when the objective is tractable filtering, simulation,
and control. It also supports the idea that exact joint laws are often replaced
by modular approximations that must then be validated.

\paragraph{What remains bank-specific.}
The exact factorization used in the proposal remains a proposal-specific
approximation. The literature cannot prove that the bank's chosen decomposition
is exact rather than a pragmatic modeling choice.

\section{Belief approximation and compression}

\paragraph{Gap in current proposal.}
Once the proposal introduces latent state and belief-state control, it inherits a
tractability problem: posterior beliefs can be very high-dimensional. A committee
will ask how the bank will retain a usable information state without pretending
that a low-dimensional approximation is exact.

\paragraph{Where the gap appears.}
This issue appears in the state-space assembly section's belief-state equations
and filtering discussion, especially where the proposal suggests finite-
dimensional posterior summaries.

\paragraph{Literatures that can close or tighten it.}
The Bayesian filtering and POMDP literatures provide the formal backdrop. In
addition, belief-approximation and finite-memory references can tighten the
story. The POMDP approximation literature
\citep{smallwood1973pomdp,kaelbling1998pomdp} and more explicit aggregation /
finite-memory references such as \citet{bean1987aggregation} can support the
idea that reduced belief summaries may be legitimate if their approximation role
is made explicit. The proposal's new filtering section already cites the core
state-space and sequential Monte Carlo references
\citep{durbinkoopman2012ssm,westharrison1997bfdm,doucet2001smc}.

\paragraph{What the literature actually supports.}
The literature supports the distinction between the exact sufficient statistic
(the full posterior distribution) and an operational approximation used for
tractability. It also supports the view that approximate filtering and
finite-dimensional compression are legitimate implementation devices.

\paragraph{What remains bank-specific.}
The exact compressed belief state the bank should use is not identified by the
literature. Its quality must be evaluated in the bank's setting by calibration,
policy-value stability, and sensitivity analysis.

\section{Incremental baseline semantics in dynamic settings}

\paragraph{Gap in current proposal.}
The proposal is strong on incremental valuation, but dynamic settings make the
baseline more subtle. Value depends not only on the current action but also on
the downstream policy bundle, the horizon, and the counterfactual path.

\paragraph{Where the gap appears.}
This gap appears in proposal objective (eq:proposal-objective),
NPV functional (eq:panel-npv-functional), component delta-NPV
(eq:component-delta-npv), dynamic-policy system (eq:dynamic-policy-system), and
policy value (eq:policy-value). It also appears in the new Bellman-style section
of the proposal.

\paragraph{Literatures that can close or tighten it.}
Causal inference and dynamic treatment-regime literatures are most helpful here
\citep{murphy2003dynamic,imbensrubin2015causal,hernanrobins2020causal,
puterman1994mdp}. They support the distinction between:
- the value of a current action under a fixed downstream policy,
- and the value of a newly optimized policy.

\paragraph{What the literature actually supports.}
These references strongly support the requirement that baseline action,
downstream policy, horizon, and information set must be explicit if a dynamic
incremental value object is to be interpretable.

\paragraph{What remains bank-specific.}
The bank must still choose the approved baseline semantics, the declared
policy-bundle convention, and the exact bridge from one-step incremental cash
flow to long-horizon value.

\section{Wallet and latent demand identification}

\paragraph{Gap in current proposal.}
This is one of the hardest remaining modeling gaps. The proposal's cross-sell
and cannibalization sections rely on outside wallet, movable wallet, competitor
substitution, and promotion-timing objects that are central to bank value but
only partially observed.

\paragraph{Where the gap appears.}
The gap is concentrated in wallet accounting (eq:wallet-accounting),
spend-source panel (eq:spend-source-panel), spend-source decomposition
(eq:spend-source-decomposition), movable-wallet object
(eq:movable-card-wallet), incremental spend-source
(eq:incremental-spend-source), new-card lift (eq:new-card-lift), bank-spend
lift (eq:bank-spend-lift), cannibalization (eq:cannibalization), and
cannibalization-to-NPV (eq:cannibalization-to-npv).

\paragraph{Literatures that can close or tighten it.}
Payment-choice and instrument-use models already in the bibliography are a first
anchor \citep{koulayev2012payment,prelec2001creditcard,agarwal2010rewards}. The
causal/uplift literature helps with treatment-effect framing
\citep{rosenbaum1983propensity,lo2002true,wager2018causalforests,
chernozhukov2018dml}. The new state-space and belief-state references justify
modeling outside wallet and movable wallet as latent state rather than observed
fact \citep{smallwood1973pomdp,kaelbling1998pomdp,
durbinkoopman2012ssm,doucet2001smc}. Banking-specific dynamic control papers
such as \citet{trench2003bankone} and \citet{thomas2011creditcardmdp} add
weight that a dynamic value system can legitimately include credit-card profit
and customer-state transitions even when not all drivers are observed directly.

\paragraph{What the literature actually supports.}
These literatures support separating adoption, use, routing, spend source, and
policy response. They also support treating wallet-share and substitution objects
as latent or partially anchored state components instead of pretending they are
directly observed.

\paragraph{What remains bank-specific.}
No public literature identifies this issuer's exact competitor-wallet capture,
movable-wallet fraction, or promotion-decay profile. Those remain bank-specific
latent demand objects to be estimated, bounded, or scenario-tagged with internal
anchors.

\section{Relationship value integration}

\paragraph{Gap in current proposal.}
The proposal already treats relationship spillovers cautiously, but the dynamic
model raises a harder question: when should relationship value be part of the
core NPV object, and when should it remain a separately reported spillover?

\paragraph{Where the gap appears.}
This matters at relationship value (eq:relationship-value), in the
valuation-semantics section, and in the allowed-use and evidence-grade
sections.

\paragraph{Literatures that can close or tighten it.}
The current CLV and customer-equity references
\citep{berger1998clv,rust2004return,venkatesan2004clv,bolton1998duration}
support the general idea that multi-period customer value can include multiple
revenue and cost channels. Banking-specific relationship and value references
such as \citet{lewis2005relationshippricing} and \citet{egan2022bankvalue} can
help explain why relationship-level value and deposit franchise value matter in
bank economics, even if they do not directly identify card-driven spillovers.

\paragraph{What the literature actually supports.}
The literature supports reporting relationship value separately or giving it an
explicit evidence grade. It supports the idea that relationship economics matter
for bank value and customer policy design.

\paragraph{What remains bank-specific.}
The literature does not justify casually adding relationship spillover to core
causal card NPV. Whether and when to include it remains an issuer-specific
identification and governance choice.

\section{Boundary between fixed-policy valuation and optimization}

\paragraph{Gap in current proposal.}
The proposal now contains dynamic-policy, belief-state, and Bellman language.
That strengthens the theoretical backbone, but it also creates a boundary risk:
is the proposal still a valuation component under a declared policy, or is it now
quietly becoming an optimizer?

\paragraph{Where the gap appears.}
This gap appears in dynamic-policy system (eq:dynamic-policy-system),
policy value (eq:policy-value), the state-space assembly section, the
allowed-use taxonomy, and the governance / migration sections.

\paragraph{Literatures that can close or tighten it.}
The MDP and dynamic treatment-regime references support fixed-policy value and
later optimization \citep{puterman1994mdp,murphy2003dynamic}. The off-policy
evaluation and safe-improvement literature
\citep{jiangli2016dr,huwager2021pomdpope,laroche2019spibb,
khraishiokhrati2022offlinecredit} supports an offline-first ladder:
- fixed-policy simulation,
- challenger offline evaluation,
- and only then constrained policy improvement.

\paragraph{What the literature actually supports.}
These literatures support the proposal's current direction if it says clearly:
fixed-policy valuation is inside current scope; optimization is a future,
governed extension. They support a staged path, not immediate production RL.

\paragraph{What remains bank-specific.}
The bank must still determine logging adequacy, overlap/support quality, model
risk appetite, safe rollout threshold, and whether a valuation component should
ever directly feed an optimizing policy layer in production.

\section{What the literature can close more fully than bank evidence can}

Some gaps are especially amenable to literature support. In particular, the
proposal can become much stronger theoretically by citing literatures on:
\begin{enumerate}[label=\arabic*.]
  \item explicit filtration and belief-state sufficiency,
  \item dynamic state aggregation and approximation,
  \item common-shock macro factorization,
  \item the difference between fixed-policy value and optimized-policy value,
  \item and conservative offline policy evaluation / safe improvement.
\end{enumerate}
These are places where the literature can do real work for the proposal by
clarifying formal language and valid abstraction.

By contrast, some gaps remain stubbornly empirical even after the literature is
added:
\begin{enumerate}[label=\arabic*.]
  \item issuer-specific outside-wallet capture,
  \item movable-wallet magnitude,
  \item relationship spillover causality,
  \item exact behavior-policy logging quality,
  \item and promotion thresholds for live challenger deployment.
\end{enumerate}
Those should remain explicitly flagged as bank-specific.

\section{Writing and citation discipline for the proposal and companion notes}

This note recommends the following discipline for the main proposal and the
stochastic-process note.
\begin{enumerate}[label=\arabic*.]
  \item If a statement is about why a modeling language is appropriate, cite the
        formal literature.
  \item If a statement is about a dynamic equation already in the proposal,
        cross-reference the proposal equation directly.
  \item If a statement is about a new assembly equation, label it as proposal
        derivation or assembly derivation.
  \item If a statement is about an issuer-specific quantity, say that explicitly
        and do not let the literature citation imply otherwise.
  \item If a claim is only scenario, benchmark, or prior support, label it as
        such and do not let it masquerade as identification.
\end{enumerate}

The note should therefore be read as a committee-defense crosswalk, not as a
claim that the literatures close every gap completely.

\section{Residual bank-specific open items}

Even after the literature crosswalk, the bank still has to resolve:
\begin{enumerate}[label=\arabic*.]
  \item the exact dynamic state unit and aggregation level by use case,
  \item the actual observed feature set and source-system timing contract,
  \item the latent-state approximation used in production,
  \item the internal anchors for wallet and substitution objects,
  \item the relationship-value inclusion rule,
  \item the approved valuation-semantics bundle,
  \item the behavior-policy logging and overlap diagnostics for offline
        evaluation,
  \item and the governance threshold separating fixed-policy valuation from
        approved challenger optimization.
\end{enumerate}

These residuals are not a weakness of the note. They are the correct boundary
between literature-supported theory and issuer-specific implementation.

\section{Conclusion}

Yes, there are real literatures that can plug several of the remaining gaps in
the current proposal. They are strongest where the proposal needs:
- formal dynamic-language support,
- explicit filtration and belief-state support,
- structured aggregation and common-shock state treatment,
- and disciplined boundaries between valuation, policy evaluation, and later
  optimization.

Those literatures can materially strengthen the proposal. But they do not remove
the need for issuer-specific evidence on latent wallet capture, relationship
spillovers, dynamic support, and deployment thresholds. The correct defense is
therefore not that the literature solves the whole bank problem. It is that the
literature sharply improves the theoretical and modeling discipline of the
proposal while leaving the genuinely issuer-specific quantities where they
belong: in internal data, experiments, quasi-experiments, and governance.

\bibliographystyle{plainnat}
\bibliography{credit_card_npv_component_refs}

\end{document}
